% reference_columns.tex — opt-in flowing two/three-column reference layout.
% Load after preamble.tex. Font sizes are explicit layout declarations, never
% inferred from content length. Ordinary multicol balancing/warnings remain live.
% Declare body font size/baseline before this file. For an authored topic break,
% end the paragraph before the native \columnbreak (\par\columnbreak); a
% columnbreak inserted into a paragraph marks its line end, not a word boundary.
% <hllpdf:architecture-guard>
\ifdefined\hllpdfcoreloaded\else
  \errmessage{hllpdf: reference_columns.tex requires preamble.tex first}
  \expandafter\endinput
\fi
\ifdefined\hllpdfreferencecolumnsloaded
  \errmessage{hllpdf: reference_columns.tex loaded more than once}
  \expandafter\endinput
\fi
\def\hllpdfreferencecolumnsloaded{1}
% </hllpdf:architecture-guard>
\usepackage{multicol}
\providecommand{\ReferenceColumnCount}{2}
\providecommand{\ReferenceColumnGap}{5mm}
\providecommand{\ReferenceColumnRule}{.3pt}
\providecommand{\ReferenceColumnSkip}{6pt}
\providecommand{\ReferenceBodyFontSize}{9}
\providecommand{\ReferenceBodyBaseline}{11}
\providecommand{\ReferenceMathFontSize}{\ReferenceBodyFontSize}
\makeatletter
% Authored sizes are independent of the prose baseline. TeX-distribution
% Computer Modern outlines are genuinely scalable: retain the standard
% ten-point design for main mathematics and the seven/five-point optical
% designs for the declared script tiers. Scaling is uniform font scaling;
% no PDF box, matrix, glyph outline or direction primitive is altered.
\@tempdima=\dimexpr\ReferenceMathFontSize pt*7/10\relax
\edef\hllpdf@referenceusualscript{\strip@pt\@tempdima}
\@tempdima=\dimexpr\ReferenceMathFontSize pt/2\relax
\edef\hllpdf@referenceusualscriptscript{\strip@pt\@tempdima}
\providecommand{\ReferenceMathScriptSize}{\hllpdf@referenceusualscript}
\providecommand{\ReferenceMathScriptScriptSize}{\hllpdf@referenceusualscriptscript}
\newcommand{\hllpdf@referencecheckmathsizes}{%
  \ifdim\ReferenceMathScriptScriptSize pt>\z@\else
    \PackageError{hebrew-english-latex-studio}{Reference math sizes must be positive}%
      {Declare positive main, script and scriptscript sizes before reference_columns.tex.}%
  \fi
  \ifdim\ReferenceMathScriptSize pt<\ReferenceMathScriptScriptSize pt
    \PackageError{hebrew-english-latex-studio}{Reference math script sizes must be ordered}%
      {The script size must be at least the scriptscript size.}%
  \fi
  \ifdim\ReferenceMathFontSize pt<\ReferenceMathScriptSize pt
    \PackageError{hebrew-english-latex-studio}{Reference math sizes must be ordered}%
      {The main math size must be at least the script size.}%
  \fi}
\hllpdf@referencecheckmathsizes
% Each configured tier receives its intended optical design. At larger local
% heading sizes these same deterministic ranges naturally select a larger
% design. These named families never replace the core's normal/bold families.
\@tempdima=\dimexpr(\ReferenceMathFontSize pt+\ReferenceMathScriptSize pt)/2\relax
\edef\hllpdf@referencemathupper{\strip@pt\@tempdima}
\@tempdima=\dimexpr(\ReferenceMathScriptSize pt+\ReferenceMathScriptScriptSize pt)/2\relax
\edef\hllpdf@referencemathlower{\strip@pt\@tempdima}
\newcommand{\hllpdf@referencescalableshape}[7]{%
  \edef\hllpdf@referencefaces{%
    <-\hllpdf@referencemathlower>#5 %
    <\hllpdf@referencemathlower-\hllpdf@referencemathupper>#6 %
    <\hllpdf@referencemathupper->#7}%
  \DeclareFontShape{#1}{#2}{#3}{#4}{\hllpdf@referencefaces}{}}
\DeclareFontFamily{OT1}{hllrefr}{}
\hllpdf@referencescalableshape{OT1}{hllrefr}{m}{n}{cmr5}{cmr7}{cmr10}
\hllpdf@referencescalableshape{OT1}{hllrefr}{bx}{n}{cmbx5}{cmbx7}{cmbx10}
\DeclareFontShape{OT1}{hllrefr}{m}{it}{<->cmti10}{}
\DeclareFontShape{OT1}{hllrefr}{bx}{it}{<->cmbxti10}{}
\DeclareFontFamily{OML}{hllrefmi}{\skewchar\font127}
\hllpdf@referencescalableshape{OML}{hllrefmi}{m}{it}{cmmi5}{cmmi7}{cmmi10}
\hllpdf@referencescalableshape{OML}{hllrefmi}{b}{it}{cmmib5}{cmmib7}{cmmib10}
\DeclareFontFamily{OMS}{hllrefsy}{\skewchar\font48}
\hllpdf@referencescalableshape{OMS}{hllrefsy}{m}{n}{cmsy5}{cmsy7}{cmsy10}
\hllpdf@referencescalableshape{OMS}{hllrefsy}{b}{n}{cmbsy5}{cmbsy7}{cmbsy10}
\DeclareFontFamily{OMX}{hllrefex}{}
\DeclareFontShape{OMX}{hllrefex}{m}{n}{<->cmex10}{}
\DeclareFontFamily{OT1}{hllrefsf}{}
\DeclareFontShape{OT1}{hllrefsf}{m}{n}{<->cmss10}{}
\DeclareFontShape{OT1}{hllrefsf}{bx}{n}{<->cmssbx10}{}
\DeclareFontFamily{OT1}{hllreftt}{}
\DeclareFontShape{OT1}{hllreftt}{m}{n}{<->cmtt10}{}
\DeclareFontFamily{U}{hllrefmsa}{}
\hllpdf@referencescalableshape{U}{hllrefmsa}{m}{n}{msam5}{msam7}{msam10}
\DeclareFontFamily{U}{hllrefmsb}{}
\hllpdf@referencescalableshape{U}{hllrefmsb}{m}{n}{msbm5}{msbm7}{msbm10}
\DeclareMathVersion{hllreference}
\DeclareMathVersion{hllreferencebold}
\SetSymbolFont{operators}{hllreference}{OT1}{hllrefr}{m}{n}
\SetSymbolFont{letters}{hllreference}{OML}{hllrefmi}{m}{it}
\SetSymbolFont{symbols}{hllreference}{OMS}{hllrefsy}{m}{n}
\SetSymbolFont{largesymbols}{hllreference}{OMX}{hllrefex}{m}{n}
\SetSymbolFont{operators}{hllreferencebold}{OT1}{hllrefr}{bx}{n}
\SetSymbolFont{letters}{hllreferencebold}{OML}{hllrefmi}{b}{it}
\SetSymbolFont{symbols}{hllreferencebold}{OMS}{hllrefsy}{b}{n}
\SetSymbolFont{largesymbols}{hllreferencebold}{OMX}{hllrefex}{m}{n}
\newcommand{\hllpdf@referencealphabets}[1]{%
  \SetSymbolFont{AMSa}{#1}{U}{hllrefmsa}{m}{n}%
  \SetSymbolFont{AMSb}{#1}{U}{hllrefmsb}{m}{n}%
  \SetMathAlphabet{\mathbf}{#1}{OT1}{hllrefr}{bx}{n}%
  \SetMathAlphabet{\mathsf}{#1}{OT1}{hllrefsf}{m}{n}%
  \SetMathAlphabet{\mathit}{#1}{OT1}{hllrefr}{m}{it}%
  \SetMathAlphabet{\mathtt}{#1}{OT1}{hllreftt}{m}{n}}
\forcsvlist{\hllpdf@referencealphabets}{hllreference,hllreferencebold}
\SetMathAlphabet{\mathsf}{hllreferencebold}{OT1}{hllrefsf}{bx}{n}
\SetMathAlphabet{\mathit}{hllreferencebold}{OT1}{hllrefr}{bx}{it}
\def\hllpdf@referencenormalversion{normal}
\def\hllpdf@referenceboldversion{bold}
\def\hllpdf@referencemathversion{hllreference}
\def\hllpdf@referenceboldmathversion{hllreferencebold}
% NFSS math-size declarations are otherwise global and shared by versions.
% Install only a LOCAL S@<selected text size> binding, then let native NFSS
% load the configured version. Ratios follow actual local sizes, so heading
% mathematics and amsmath's smaller text use the same authored hierarchy.
\newcommand{\hllpdf@referencelocalmathsizes}{%
  \ifcsdef{hllpdf@referenceoriginalS@\f@size @captured}{}{%
    \csdef{hllpdf@referenceoriginalS@\f@size @captured}{1}%
    \csletcs{hllpdf@referenceoriginalS@\f@size}{S@\f@size}}%
  \dimen@=\f@size pt\relax
  \dimen@ii=\ReferenceBodyFontSize pt\relax
  \@tempdima=\dimexpr\ReferenceMathFontSize pt*\number\dimen@/\number\dimen@ii\relax
  \edef\hllpdf@referenceactualmath{\strip@pt\@tempdima}%
  \@tempdima=\dimexpr\ReferenceMathScriptSize pt*\number\dimen@/\number\dimen@ii\relax
  \edef\hllpdf@referenceactualscript{\strip@pt\@tempdima}%
  \@tempdima=\dimexpr\ReferenceMathScriptScriptSize pt*\number\dimen@/\number\dimen@ii\relax
  \edef\hllpdf@referenceactualscriptscript{\strip@pt\@tempdima}%
  \expandafter\edef\csname S@\f@size\endcsname{%
    \gdef\noexpand\tf@size{\hllpdf@referenceactualmath}%
    \gdef\noexpand\sf@size{\hllpdf@referenceactualscript}%
    \gdef\noexpand\ssf@size{\hllpdf@referenceactualscriptscript}%
    \noexpand\math@fontstrue}}
\newif\ifhllpdf@referencemathactive
\newcommand{\hllpdf@referenceactivatemath}{%
  \ifhllpdf@referencemathactive\else
    \hllpdf@referencemathactivetrue
    \let\hllpdf@referenceoriginalmathversion\mathversion
    % Native boldmath, unboldmath and amsmath's boldsymbol request the normal
    % and bold names. Route only these requests to the matching reference
    % versions; independently authored custom versions pass through intact.
    \protected\def\mathversion##1{%
      \edef\hllpdf@referencerequestedversion{##1}%
      \ifx\hllpdf@referencerequestedversion\hllpdf@referencenormalversion
        \hllpdf@referenceoriginalmathversion{hllreference}%
      \else\ifx\hllpdf@referencerequestedversion\hllpdf@referenceboldversion
        \hllpdf@referenceoriginalmathversion{hllreferencebold}%
      \else
        \hllpdf@referenceoriginalmathversion{##1}%
      \fi\fi}%
    \let\hllpdf@referenceoriginalmathsettings\glb@settings
    \def\glb@settings{%
      \ifx\math@version\hllpdf@referencemathversion
        \hllpdf@referencelocalmathsizes
      \else\ifx\math@version\hllpdf@referenceboldmathversion
        \hllpdf@referencelocalmathsizes
      \else
        \ifcsdef{hllpdf@referenceoriginalS@\f@size @captured}{%
          \csletcs{S@\f@size}{hllpdf@referenceoriginalS@\f@size}}{}%
      \fi\fi
      \hllpdf@referenceoriginalmathsettings}%
    \mathversion{\math@version}%
  \fi}
\newif\ifhllpdf@referencecolumnsactive
\newif\ifhllpdf@referencewideactive
% Native sectioning keeps its own font, rule, counter, star and anchor policy.
% Only active grouped reference layouts add natural reading-edge heading lines
% and suppress automatic word hyphenation. In LuaTeX \raggedright follows the
% paragraph direction: physical left in English and physical right in Hebrew.
% Prepend setup to style argument 6: a theme may end that style with a macro
% which consumes the heading title, so appending tokens would change its input.
\newif\ifhllpdf@referenceheadingsactive
\newcommand{\hllpdf@referenceheadingsetup}{%
  \raggedright
  \hyphenpenalty=10000\relax}
\newcommand{\hllpdf@referenceactivateheadings}{%
  \ifhllpdf@referenceheadingsactive\else
    \hllpdf@referenceheadingsactivetrue
    \let\hllpdf@referenceoriginalstartsection\@startsection
    \def\@startsection##1##2##3##4##5##6{%
      \hllpdf@referenceoriginalstartsection{##1}{##2}{##3}{##4}{##5}%
        {\hllpdf@referenceheadingsetup ##6}}%
  \fi}
% Presentation hooks run after the declared local font has been selected, so
% a theme can derive display spacing from the actual reference baseline.
% An environment's early begin hook would still see the outer article size.
% This private empty hook preserves standalone layout and either load order.
\providecommand{\hllpdf@referencebodysetup}{}
\newcommand{\ReferenceSelectBodyFont}{%
  \ifdim\ReferenceBodyFontSize pt>\z@\else
    \PackageError{hebrew-english-latex-studio}{Reference body font size must be positive}%
      {Use an explicit positive point size, for example 9.}%
  \fi
  \ifdim\ReferenceBodyBaseline pt<\ReferenceBodyFontSize pt
    \PackageError{hebrew-english-latex-studio}{Reference baseline must not be smaller than the font size}%
      {Declare a readable baseline in points, for example 11 for a 9-point font.}%
  \fi
  \hllpdf@referencecheckmathsizes
  \fontsize{\ReferenceBodyFontSize}{\ReferenceBodyBaseline}\selectfont
  % Compact entries start at their local reading edge under every theme.
  % These assignments belong to the grouped layout, not to raggedright in a
  % presentation overlay; the surrounding document's paragraph policy returns.
  \setlength{\parindent}{0pt}
  \setlength{\parskip}{0pt plus .5pt}
  \ifhllpdf@referencecolumnsactive
    \hllpdf@referenceactivateheadings
    \hllpdf@referenceactivatemath
  \else\ifhllpdf@referencewideactive
    \hllpdf@referenceactivateheadings
    \hllpdf@referenceactivatemath
  \fi\fi
  \hllpdf@referencebodysetup}
\newcommand{\hllpdf@referenceoutercheck}[1]{%
  \ifinner
    \PackageError{hebrew-english-latex-studio}{#1 requires outer text mode}%
      {Do not place this environment in a box, minipage or float.}%
  \fi
  \if@twocolumn
    \PackageError{hebrew-english-latex-studio}{#1 requires a one-column document class}%
      {Use referencecolumns for the local column layout; omit the class twocolumn option.}%
  \fi
  \ifnum\@listdepth>\z@
    \PackageError{hebrew-english-latex-studio}{#1 cannot be nested in a list}%
      {Close the list before starting this layout environment.}%
  \fi}
\newcommand{\hllpdf@referenceforbidden}[1]{%
  \PackageError{hebrew-english-latex-studio}{#1 is not supported inside referencecolumns}%
    {Use a nonfloating local-width item, or close the columns and use referencewide.}}
\newenvironment{referencecolumns}[1][\ReferenceColumnCount]{%
  \par
  \hllpdf@referenceoutercheck{referencecolumns}%
  \ifhllpdf@referencecolumnsactive
    \PackageError{hebrew-english-latex-studio}{referencecolumns cannot be nested}%
      {Close the existing columns before starting a new set.}%
  \fi
  \ifhllpdf@referencewideactive
    \PackageError{hebrew-english-latex-studio}{referencecolumns cannot be nested inside referencewide}%
      {Close the full-width block before starting columns.}%
  \fi
  \edef\hllpdf@referencecount{#1}%
  \def\hllpdf@referencetwo{2}\def\hllpdf@referencethree{3}%
  \ifx\hllpdf@referencecount\hllpdf@referencetwo\else
    \ifx\hllpdf@referencecount\hllpdf@referencethree\else
      \PackageError{hebrew-english-latex-studio}{referencecolumns accepts only 2 or 3 columns}%
        {Set ReferenceColumnCount or the environment's optional argument to 2 or 3.}%
    \fi
  \fi
  \setlength{\columnsep}{\ReferenceColumnGap}%
  \ifdim\columnsep>\z@\else
    \PackageError{hebrew-english-latex-studio}{Reference column gap must be positive}%
      {Use an explicit positive length, for example 5mm.}%
  \fi
  \setlength{\multicolsep}{\ReferenceColumnSkip}%
  \ifdim\multicolsep<\z@
    \PackageError{hebrew-english-latex-studio}{Reference column skip must not be negative}%
      {Use zero or a positive length.}%
  \fi
  \setlength{\columnseprule}{\ReferenceColumnRule}%
  \ifdim\columnseprule<\z@
    \PackageError{hebrew-english-latex-studio}{Reference column rule must not be negative}%
      {Use zero to omit separators or an explicit positive length.}%
  \fi
  \hllpdf@referencecolumnsactivetrue
  \ReferenceSelectBodyFont
  % LuaLaTeX's RTL page direction already places LR@column@boxes from right
  % to left. RLmulticolcolumns would reverse that order a second time. Keep
  % the native arrangement, and normalize physical assembly to the PAGE
  % direction, including when the final paragraph is a minority-language one.
  % Babel normalizes the regular multicol output routine, but not page@sofar's
  % direct final balanced assembly. This scoped patch covers both paths.
  \LRmulticolcolumns
  \pretocmd{\page@sofar}{\bodydir\pagedir}{}{%
    \PackageError{hebrew-english-latex-studio}{Cannot normalize reference column output direction}%
      {Unsupported multicol output routine.}}%
  \def\@float##1{\hllpdf@referenceforbidden{Floats}}%
  \def\@dblfloat##1{\hllpdf@referenceforbidden{Floats}}%
  \def\marginpar{\hllpdf@referenceforbidden{Margin notes}}%
  \ifdefined\longtable
    \def\longtable{\hllpdf@referenceforbidden{longtable}}%
  \fi
  \ifdefined\flowtable
    \def\flowtable{\hllpdf@referenceforbidden{flowtable}}%
  \fi
  % Native multicol permits up to 12pt of final balancing overflow. For a
  % reference page the authored atomic frames must fit the real text area;
  % otherwise use multicol's ordinary next-page path. Native legal glue shrink
  % remains available when forming balanced columns within that real area.
  \maxbalancingoverflow=0pt\relax
  % multicol normally packs each final trial into the clamped target height
  % BEFORE discovering that the trial is too tall. That emits an Overfull vbox
  % even when its native fallback subsequently reflows perfectly. Measure the
  % natural box first and reject a column exceeding the AVAILABLE PAGE HEIGHT
  % without constructing it. Do not compare it to the smaller equal-balancing
  % target: legal glue shrink may fit that target, and prematurely rejecting it
  % would turn a valid three-column balance into full-height two-column flow.
  % The native vfuzz=0 diagnostic policy remains active in BOTH branches; no
  % warning threshold, checker budget or ordinary output routine is relaxed.
  \patchcmd{\balance@columns}{\vfuzz\z@}{%
    \vfuzz\z@
    \setbox\@tempboxa\vbox{\unvcopy\@tempcnta}%
    \ifdim\ht\@tempboxa>\@colroom
      \too@badtrue
    \else}{}{%
      \PackageError{hebrew-english-latex-studio}%
        {Cannot preflight reference column balancing}%
        {The installed multicol final-column packing routine is unsupported.}}%
  \patchcmd{\balance@columns}{\ifnum\badness>\@M}{%
    \fi\ifnum\badness>\@M}{}{%
      \PackageError{hebrew-english-latex-studio}%
        {Cannot finish reference column balancing preflight}%
        {The installed multicol final-column badness test is unsupported.}}%
  \begin{multicols}{\hllpdf@referencecount}%
  \raggedcolumns
  \ignorespaces
}{%
  \par\end{multicols}%
  \ignorespacesafterend}
% A wide block is an explicit transition between column runs. It does not
% silently interrupt/restart multicol and does not resize any supplied content.
\newenvironment{referencewide}{%
  \par
  \ifhllpdf@referencecolumnsactive
    \PackageError{hebrew-english-latex-studio}{referencewide cannot be used inside referencecolumns}%
      {Close referencecolumns, write the wide block, then start a new column run.}%
  \fi
  \ifhllpdf@referencewideactive
    \PackageError{hebrew-english-latex-studio}{referencewide cannot be nested}%
      {Close the existing wide block first.}%
  \fi
  \hllpdf@referenceoutercheck{referencewide}%
  \hllpdf@referencewideactivetrue
  \ReferenceSelectBodyFont
  \addvspace{\ReferenceColumnSkip}%
  \ignorespaces
}{%
  \par\addvspace{\ReferenceColumnSkip}%
  \ignorespacesafterend}
\makeatother
